\chapter{Wykład II}
\section{Pojęcie podgrupy i przykłady. }
\section{Rząd elementu grupy.}
\section{Generatory grupy, grupy skończenie generowane.}
\section{Grupa dihedralna $D_{n}$}

\begin{context}[Notacja]
	W tej notatce $D_{n}$ oznacza \textbf{rzędu 2n}, reprezentującą symetrie 2n-kąta foremnego.
\end{context}
\begin{enumerate}[label=\arabic*.]
	\item \textbf{Definicja geometryczna} 
		
		Grupa dihedralna(diedralna) $D_{n}\left( n\ge 3 \right)$ jest \textbf{grupą symetrii n-kąta foremnego}. Składa się ona ze wszystkich izometrii płaszczyzny(obrotów i symetrii osiowych), które przeprowadzają n-kąt foremny na samego siebie.
	\item \textbf{Elementy i rząd}

		\begin{itemize}
			\item \textbf{Fakt:} Grupa $D_{n}$ ma dokładnie 2n elementów. Składa się z n obrotów i n symetrii.  
		\end{itemize}
\end{enumerate}

